\documentclass[10pt,a4paper]{amsart}
\usepackage[margin=19mm]{geometry}
\usepackage{amsmath,amssymb,mathtools}
\usepackage[scr=rsfs,cal=euler]{mathalfa}
\usepackage{xfrac}
\usepackage[unicode,hidelinks,bookmarks=false]{hyperref}
\newcommand{\ie}{i.e.\ }
\newcommand{\cf}{cf.\ }
\newcommand{\resp}{respectively\ }
\newcommand{\Subsection}{\subsection}
\providecommand{\nobreakdash}{\nobreak\hbox{-}}
\setlength{\emergencystretch}{2em}
\multlinegap0pt
\usepackage{tikz}
\usetikzlibrary{cd}
\usepackage{mathtools}
\usepackage{amssymb}
\newtheorem{prop}{Proposition}
\DeclareMathOperator{\BAB}{B\otimes_{A} B}
\DeclareMathOperator{\Sp}{Spec}
\title{Real radiciality and monoreal extensions}
\author{Goulwen Fichou, Jean-Philippe Monnier, Ronan Quarez}
\date{Source: arXiv:2607.03408v1 (CC BY 4.0)}
\begin{document}
\maketitle

\noindent\emph{Source and licence.} Real radiciality and monoreal extensions. Version arXiv:2607.03408v1 (CC BY 4.0), distributed by arXiv under the Creative Commons Attribution 4.0 International licence, http://creativecommons.org/licenses/by/4.0/, as recorded in the arXiv licence metadata for this exact version and shown on the abstract page https://arxiv.org/abs/2607.03408v1. Full licence notice: \texttt{licenses/arxiv-2607.03408-CC-BY-4.0.txt}.\par\smallskip

\noindent\textit{Editorial scope.} Counted unit: the article's prop rrsaturdeux (source0.tex lines 965--972). The article's complete original proof of it is delivered with it (source0.tex lines 973--994, one \texttt{proof} environment of 22 lines). No mathematics is written, repaired or re-derived: the statement and the proof are the article's own lines, in the article's own notation, and no retained line is substituted or re-flowed.  No further source-local context is needed: every cross-reference of every retained line resolves inside the two delivered spans.  Nothing was omitted from any retained span, so the package carries no omission record.  No retained line contains a citation, so the wrapper carries no bibliography and no bibliography entry.  The retained lines contain the article's own diagram code, which the wrapper typesets with the standard tikz packages; no image file is delivered and the package declares no asset dependency.  Editorial formatting only, in addition: an AMS \texttt{amsart} preamble that restates the article's own declarations and macros used by the retained lines, namely the environments \texttt{prop} on separate counters, together with the standard packages those lines need.  The retained environments print in this document as Proposition 1 in order of appearance, whereas the article prints the same items with the numbering of its own full document.  The complete native source of the article as distributed in the arXiv e-print for this exact version is bundled unchanged as \texttt{sources/204384/source0.tex}.
\begin{prop}\label{rrsaturdeux}
Let $A\xrightarrow{i}  B$ be an extension of rings. Then, the following properties are equivalent:
\begin{enumerate}
\item $i$ is  real radicial.
%\item $\widehat{A}^r_B=B$.
\item The map $\Sp_r (\Delta):\Sp_r B\to \Sp_r (\BAB)$ is surjective where $\Delta:\BAB\to B$ maps $b_1\otimes b_2$ to $b_1b_2$ is the canonical diagonal morphism.
\end{enumerate}
\end{prop}

\begin{proof}
\textbf{(2) $\implies$ (1):}
Assume (2). Let $\psi_1:B\to R$ and $\psi_2:B\to R$ be two morphisms 
such that $\psi_1\circ  i=\psi_2\circ i$. We get a morphism $\psi:\BAB\to R$ such that $\psi(b_1\otimes_A b_2)=\psi_1(b_1)\psi_2(b_2)$. By (2), the point in $\Sp_r(\BAB)$ given by $\psi$ is the image of an element of $\Sp_r B$ by $\Sp_r(\Delta)$. So, up to replacing $R$ by a real closed extension, we get 
a morphism $\psi^{\prime}:B\to R$ such that the following diagram
$$\begin{tikzcd}
\BAB \arrow[d, "\Delta" '] \arrow[r,"\psi"] & R\\
B \arrow[ur,,"\psi'" ',] &
\end{tikzcd}$$
is commutative. Let $b\in B$. We have 
$$\psi_1(b)=\psi(b\otimes_A 1)=\psi'\circ \Delta(b\otimes_A 1)=\psi'(b)=\psi'\circ \Delta(1 \otimes_A b)=\psi(1\otimes_A b)=\psi_2(b),$$ so that $\psi_1=\psi_2$.

\textbf{(1) $\implies$ (2):}
Let $\psi:\BAB\to R$ be a morphism. Composing $\psi$ with $b\mapsto b\otimes_{A}1$ and $b\mapsto 1\otimes_{A}b$, we obtain two morphisms $\psi_i:B\to R$, $i=1,2$. Clearly, $\psi_1\circ  i=\psi_2\circ  i$: for $a\in A$ we have $\psi_1\circ i(a)=\psi( i(a)\otimes_A 1)=\psi(1\otimes_A  i(a))=\psi_2\circ  i(a)$. Remark that for $b_1,b_2$ in $B$ we have $\psi(b_1\otimes_A b_2)=\psi_1(b_1)\psi_2(b_2)$. By (1), it follows that $\psi_1=\psi_2$. We get a commutative diagram
$$\begin{tikzcd}
\BAB \arrow[d, "\Delta" '] \arrow[r,"\psi"] & R\\
B \arrow[ur,,"\psi'" ',] &
\end{tikzcd}$$
by setting $\psi'(b)=\psi_1(b)=\psi_2(b)$. We clearly get $(\Sp_r \Delta)(\psi')=\psi$, namely the point in 
$\Sp_r (\BAB)$ given by $\psi$ is the image of the point in 
$\Sp_r B$ given by $\psi'$.
\end{proof}

\end{document}
